\chapter{虚拟存储器}

\section{虚拟存储器概念}
\concept{为什么需要虚拟存储器}
\begin{points}
  \item 传统物理存储管理要求作业运行前全部装入内存，运行期间常驻内存，限制了大作业和多作业运行。
  \item 物理扩充内存会增加成本，因此需要从逻辑上扩充内存。
  \item 虚拟存储器打破“作业必须全部装入内存才能运行”的限制。
\end{points}
\plain{虚拟存储器首先解决的是“内存装不下整个作业怎么办”这个问题：不再要求程序一开始全部进内存，而是先装一部分，边运行边调入。}
\exam{简答题常从“为什么需要虚拟存储器”切入，答案核心是：突破全部装入限制、逻辑扩充内存、降低物理扩容成本。}


\concept{局部性原理}
\begin{points}
  \item 程序执行时，CPU 访问的指令和数据倾向于聚集在较小范围内。
  \item 时间局部性：刚访问过的指令或数据，短时间内可能再次访问。
  \item 空间局部性：访问某地址后，附近地址很可能被访问。
  \item 顺序局部性：程序大多数情况下顺序执行，下一条指令通常紧接当前指令。
  \item 局部性原理是虚拟存储器的理论基础。
\end{points}
\plain{程序一段时间只爱访问附近的一小块代码和数据，所以没必要一开始全装入内存。}
\exam{除了时间/空间局部性定义，还要会说它为什么能支持请求分页和工作集模型。}


\concept{虚拟存储器定义与特征}
\begin{points}
  \item 虚拟存储器是具有请求调入和置换功能、能从逻辑上扩充内存容量的存储器系统。
  \item 它是一种以时间换空间的技术：用缺页中断、磁盘换入换出换取更大的逻辑空间。
  \item 特征：离散性、多次性、对换性、虚拟性。
  \item 容量受地址结构和外存容量限制，不是无限大。
\end{points}
\plain{“虚拟”不是凭空变出内存，而是把程序访问地址和真实物理内存解耦，再靠外存做后备仓库。}
\exam{四特征是高频名词解释；回答时最好逐个落到机制上：离散对应非连续分配，多次对应分批装入，对换对应换入换出，虚拟对应逻辑扩充。}


\concept{虚拟存储器的好处与实现基础}
\begin{points}
  \item 好处：可以运行更大的程序，因为不再要求程序整体常驻内存。
  \item 好处：可以同时运行更多程序，因为每个程序通常只需把活跃页面保留在内存中。
  \item 好处：程序启动更快，因为初始装入的页面更少。
  \item 物质基础：相当数量的外存、一定容量的内存，以及能动态完成地址转换的硬件机构。
\end{points}
\plain{这部分可以连起来记：外存负责“兜底”，内存负责“放当前活跃部分”，地址变换机构负责“把逻辑访问导向当前真正所在的位置”。}
\exam{PPT 明确提了“好处”和“实现基础”，简答题不要只背定义而漏掉条件。}


\section{请求分页存储管理}
\concept{请求分页的基本思想}
\begin{points}
  \item 程序运行前只装入部分页面即可启动。
  \item 运行中访问的页面若不在内存，产生缺页中断，由 OS 调入所需页面。
  \item 若内存无空闲页框，需要选择某个页面换出。
  \item 请求分页 = 基本分页 + 请求调页 + 页面置换。
\end{points}
\plain{请求分页和基本分页的根本差别在于：基本分页默认页已经在内存，请求分页允许“访问时才发现页不在”，然后现调。}
\exam{这正是老师点名的比较点，答题时一定写出“允许页面不在内存”和“缺页中断是正常机制”这两句。}


\concept{请求分页系统的组成}
\begin{points}
  \item 请求分页建立在分页系统基础上，还需要缺页中断机构。
  \item 需要扩充页表项，加入存在位、修改位、访问位、外存地址等信息。
  \item 还需要请求调页和页面置换的软件支持。
  \item 其硬件核心仍是 MMU，它负责地址转换、合法性检查，并与页表/TLB 配合工作。
\end{points}
\plain{请求分页不是只靠一个“缺页中断”词就结束了，它还要有页表扩展和调页/置换机制一起配合。}
\exam{常考“请求分页比基本分页多了什么支持机构”。}


\concept{MMU 在请求分页中的作用}
\begin{points}
  \item MMU 接受 CPU 给出的逻辑地址，完成页号与页内位移的分解，并把逻辑地址转换成物理地址。
  \item MMU 会先结合 TLB 和页表检查该页是否在主存中，以及访问是否合法。
  \item 若页表表明该页不在内存，MMU 触发缺页异常，由操作系统调页。
  \item 因而 MMU 既负责地址转换，也承担内存保护和虚拟内存支持的重要硬件职责。
\end{points}
\plain{可以把 MMU 看成 CPU 和内存之间的“翻译兼门卫”：先翻译地址，再判断能不能进。}
\exam{PPT 对 MMU 单独成节，说明它不只是第 8 章的旧概念复读，这里要会把 MMU 和缺页、TLB、保护联系起来。}


\concept{基本分页与请求分页区别}
\begin{points}
  \item 基本分页要求进程运行所需页面已在内存中，地址转换时页表项应有效。
  \item 请求分页允许部分页面不在内存，访问未驻留页面时产生缺页中断。
  \item 因此基本分页本身不以缺页中断为常规机制；请求分页才把缺页作为正常运行过程的一部分。
\end{points}
\plain{第 8 章基本分页讲的是“怎么按页映射”；第 9 章请求分页讲的是“页不在内存时怎么办”。}
\exam{比较题回答时别只写“一个高级一个低级”，要落到“是否允许页面暂不在内存、是否需要缺页中断和置换机制”。}


\concept{页表项扩展}
\begin{points}
  \item 存在位/有效位：表示该页是否在内存。
  \item 修改位/脏位：表示页面是否被写过，换出时若为脏页需写回外存。
  \item 引用位/访问位：表示页面近期是否被访问，用于页面置换算法。
  \item 外存地址：指出页面在外存中的位置。
\end{points}
\plain{请求分页下的页表不再只是“页号到块号”的映射表，而是还要告诉 OS 这页在不在内存、最近用没用、改没改、若不在内存该去磁盘哪找。}
\exam{填空/选择最爱考存在位、修改位、访问位、外存地址各自干什么，尤其修改位决定换出时是否必须回写。}


\concept{缺页中断处理过程}
\begin{points}
  \item CPU 访问页面，发现页表项存在位为 0，产生缺页中断。
  \item 操作系统检查访问是否合法；若非法，终止进程。
  \item 若合法，查找空闲页框；若无空闲页框，执行页面置换。
  \item 若被换出页被修改过，需要写回外存。
  \item 从外存读入所需页面，更新页表和快表，重新执行引起缺页的指令。
\end{points}
\plain{这类流程题的核心顺序是：发现不在内存 -> 判断是否合法 -> 找页框/选牺牲页 -> 必要时回写 -> 调入新页 -> 改页表/TLB -> 重执行原指令。}
\exam{老师强调过请求分页地址变换，答题时不要漏掉“重新执行产生缺页的指令”，这正是它和一般中断返回到下一条指令的重要区别。}


\concept{缺页中断与一般中断的区别}
\begin{points}
  \item 缺页中断在指令执行期间发生，而一般外部中断常在一条指令执行结束后响应。
  \item 一条指令执行过程中可能涉及多个地址访问，因此可能产生多次缺页中断。
  \item 缺页中断处理完成后，要重新执行引起缺页的那条指令；一般中断返回后通常执行下一条指令。
\end{points}
\plain{缺页中断不是“这条指令做完了再说”，而是这条指令做到一半发现缺页，必须先把缺的页补齐，回来继续把这条指令做完。}
\exam{这是 PPT 明确单列的区别点，容易出选择或简答。}


\concept{请求分页中的地址变换过程}
\begin{points}
  \item CPU 给出逻辑地址后，MMU 先拆分出页号和页内位移。
  \item 先查快表 TLB；若命中且表项有效，可直接形成物理地址访问内存。
  \item 若 TLB 未命中，则进一步查页表；若页在内存，形成物理地址并可把表项装入 TLB。
  \item 若页表显示该页不在内存，则产生缺页中断，调页成功后再更新页表/TLB 并重执行指令。
\end{points}
\plain{请求分页地址变换比基本分页只多了一条分支：查到“页不在内存”时，不是直接报错，而是进入缺页处理。}
\exam{步骤题建议写成“拆地址 -> 查 TLB -> 查页表 -> 缺页则中断 -> 调页 -> 更新 -> 重执行”，过程比一句结论更重要。}


\concept{TLB 与刷新问题}
\begin{points}
  \item TLB 是页表项的高速缓存，用于加快地址转换。
  \item 进程切换时，旧进程的地址映射可能失效，因此通常需要刷新 TLB 或切换到对应地址空间标识。
  \item 请求调页或页面置换后，如果被替换页的旧页表项仍留在 TLB 中，也必须使其失效或更新。
  \item 因而 TLB 命中并不等于“永远正确”，它必须与当前页表状态保持一致。
\end{points}
\plain{TLB 快是快，但前提是里面缓存的映射没过期；进程切换和换页都会让旧缓存失真，所以要刷新。}
\exam{PPT 专门提了 TLB 刷新，说明考试可能不只考“TLB 提高速度”，还会考“什么时候必须失效/刷新”。}


\concept{缺页次数、缺页率与缺页开销}
\begin{points}
  \item 缺页次数指页面引用过程中发生缺页的总次数。
  \item 缺页率通常指缺页次数占总页面访问次数的比例。
  \item 缺页开销远高于普通内存访问，因为往往需要磁盘 I/O。
\end{points}
\plain{一次缺页不是“多查一步表”这么简单，而是可能要从磁盘把整页搬进来，代价大得多。}
\exam{常考为什么页面置换算法和页框分配策略会显著影响系统性能。}


\concept{有效访问时间 EAT}
\begin{points}
  \item 若页在主存且 TLB 命中，访问时间约为查 TLB 时间加一次内存访问时间。
  \item 若页在主存但 TLB 未命中，则通常还需访问页表并更新 TLB，开销高于命中情况。
  \item 若页不在主存，还需加上缺页中断处理与调页时间，因此开销远大于普通内存访问。
  \item EAT 本质上是按 TLB 命中率、缺页率对几种访问路径求加权平均。
\end{points}
\plain{EAT 题不要背成死公式，先把路径画清楚：命中怎么走、未命中但不缺页怎么走、缺页怎么走，再做概率加权。}
\exam{老师在第 8 章强调过 EAT，第 9 章 PPT 又把“含缺页率的 EAT”展开了，计算题要注意把缺页处理时间单独算进去。}


\concept{请求分页的优缺点}
\begin{points}
  \item 优点：页可离散存放，减少移动和紧凑开销，基本消除外部碎片。
  \item 优点：主存利用率更高，有利于提高多道程序度和系统吞吐量。
  \item 缺点：需要硬件支持和缺页中断处理，系统实现更复杂。
  \item 缺点：若缺页频繁，磁盘 I/O 和中断处理开销会明显增加。
\end{points}
\plain{请求分页的收益来自“少装、晚装、按需装”，代价来自“查得更复杂、缺页更昂贵”。}
\exam{简答题适合按“优点两条、缺点两条”组织答案，别只写优点。}


\section{页面置换算法}
\concept{最佳置换 OPT}
\begin{points}
  \item OPT 选择未来最长时间不会被访问的页面淘汰。
  \item 它缺页率最低，但需要预知未来，实际无法实现。
  \item 常作为评价其他算法的理论下界。
\end{points}
\plain{OPT 是标准答案里的“理想参照物”：真实系统做不到，但做题时常用它来说明某个访问串的最少缺页次数。}
\exam{算法题里若同时要求比较多种置换算法，OPT 往往用来作为最优基准。}


\concept{先进先出 FIFO}
\begin{points}
  \item FIFO 淘汰最早进入内存的页面。
  \item 实现简单，只需维护队列。
  \item 缺点是不考虑页面使用频率，可能淘汰仍常用页面。
  \item FIFO 可能出现 Belady 异常：页框数增加，缺页次数反而增加。
\end{points}
\plain{FIFO 只看“谁来得早”，不看“谁最近还常用”，所以很容易把热点页错踢出去。}
\exam{Belady 异常基本就是 FIFO 的招牌考点，要会解释“页框更多却缺页更多”为什么可能发生。}


\concept{最近最久未使用 LRU}
\begin{points}
  \item LRU 淘汰最近最长时间没有被访问的页面。
  \item 基于局部性原理：最近很久未用的页面，未来短期内也可能不用。
  \item 实现可用计数器、栈或近似算法，但精确 LRU 开销较高。
  \item LRU 不会出现 Belady 异常。
\end{points}
\plain{LRU 比 FIFO 更贴合程序局部性，因为它默认“刚用过的页近期大概率还会再用”。}
\exam{题目若问为什么 LRU 通常优于 FIFO，关键词就是“利用局部性原理”。}


\concept{时钟 CLOCK/二次机会算法}
\begin{points}
  \item Clock 用访问位近似 LRU，把页面组成循环队列。
  \item 指针扫描页面：若访问位为 0，淘汰；若为 1，清 0 并给第二次机会。
  \item 改进型 Clock 还考虑修改位，优先淘汰未访问且未修改的页面。
\end{points}
\plain{Clock 的本质是“用一位访问位，低成本地近似 LRU”，牺牲一点精度换来更低实现开销。}
\exam{若题目写“二次机会算法”“简单时钟算法”，都要能和访问位、循环扫描联系起来。}

\concept{页面置换题的标准做法}
\begin{points}
  \item 按页面访问串从左到右逐次处理。
  \item 每一步先判断命中还是缺页，再决定是否置换。
  \item 若算法需要维护额外状态，如 LRU 次序或 Clock 访问位，要在每一步同步更新。
  \item 答题时最好写出每一步页框内容变化，不要只报最终缺页次数。
\end{points}
\plain{这类题最怕心算跳步，一跳步就很容易后面全错。}
\exam{老师录音强调过页面置换经常是大题，过程要完整写。}


\section{页面分配与置换策略}
\concept{页面装入与清除策略}
\begin{points}
  \item 页面装入策略决定何时把页面调入主存，常见有请求式调页和预调页两类。
  \item 请求式调页按需装入，只有访问到缺页时才调入，优点是初始装入少，缺点是运行中可能频繁触发磁盘 I/O。
  \item 预调页根据局部性原理和历史访问情况，预先把若干可能很快会访问的页面装入主存，目的是减少后续缺页次数。
  \item 页面清除策略决定何时把修改过的页面写回外存：请求式清除在页面被选中换出时才写回；预约式清除则提前成批写回脏页。
\end{points}
\plain{这里本质是在回答两个“何时”问题：何时把页装进来，何时把脏页写出去。请求式更省眼前工作，预调页和预约式清除则是想用“提前做一点”换取后续更少的阻塞。}
\exam{简答题若问页面管理策略，别只写置换算法，还要能区分“装入策略”和“清除策略”，并分别说出请求式/预调页、请求式清除/预约式清除。}

\concept{从何处调入页面}
\begin{points}
  \item 文件区用于存放普通文件，通常采用离散分配；对换区用于存放对换页面，通常采用连续分配，因此对换区 I/O 一般更快。
  \item 若系统有足够对换区，可把与进程有关的页面预先复制到对换区，缺页时直接从对换区调入。
  \item 若对换区不足，则不会被修改的页面可直接从文件区调入；可能被修改的页面换出后再写入对换区，后续优先从对换区调回。
  \item UNIX 方式中，未运行过的页面通常从文件区调入；曾运行后被换出的页面，下次一般从对换区调入；若共享页面已在内存，则可直接复用。
\end{points}
\plain{做题时别把“文件区”和“对换区”当成同一种外存位置。文件区更像原始来源地，对换区更像专门为换页准备的中转仓库。}
\exam{若考“缺页时一定从对换区调页吗”，答案不是绝对的，要看系统对换区是否充足、页面是否可修改，以及是否采用 UNIX 式处理。}

\concept{最小物理块数}
\begin{points}
  \item 最小物理块数是保证进程能够正常运行所需的最少页框数，少于该值时进程无法执行。
  \item 若是单地址指令且采用直接寻址，通常至少需要 2 个物理块。
  \item 若是单地址指令且允许间接寻址，通常至少需要 3 个物理块。
  \item 若机器较复杂，指令、源地址和目标地址都可能跨页，则所需最小物理块数会更大。
\end{points}
\plain{这个概念不是在问“运行得快不快”，而是在问“最少得有几块内存，指令和操作数访问才不至于根本做不下去”。}
\exam{选择题喜欢把“最小物理块数”和“分配越多越好”混在一起。前者是能否运行的下限，后者才是性能优化问题。}

\concept{驻留集与页面分配}
\begin{points}
  \item 驻留集指某进程当前驻留内存的页面集合。
  \item 固定分配：进程运行期间分得的页框数固定。
  \item 可变分配：根据缺页率、工作集等动态调整页框数。
\end{points}
\plain{驻留集就是“这个进程当前手里实际占着的页框集合”。}
\exam{常考固定/可变分配、局部/全局置换的组合关系。}


\concept{局部置换与全局置换}
\begin{points}
  \item 局部置换：缺页时只从本进程驻留集中选择页面淘汰。
  \item 全局置换：可从全系统所有可换页面中选择淘汰，进程页框数会动态变化。
  \item 固定分配只能配合局部置换；可变分配可配合局部置换或全局置换。
\end{points}
\plain{区分局部/全局置换时，关键看“缺页时能不能拿别的进程的页框来顶”。}
\exam{这类题是概念题，不是置换算法题，别答成 FIFO/LRU。}


\concept{不能组合的策略}
\begin{points}
  \item 固定分配 + 全局置换通常不能组合。
  \item 原因：固定分配要求每个进程页框数不变，而全局置换可能把其他进程页框拿走或让本进程获得额外页框，改变驻留集大小。
\end{points}
\plain{一个要求“页框数固定不变”，另一个却允许“从别人那儿拿页框”，两者自然冲突。}
\exam{常考为什么固定分配不能配全局置换。}


\section{工作集、抖动与缺页率}
\concept{工作集模型}
\begin{points}
  \item 工作集是在某段时间窗口内进程实际访问过的页面集合。
  \item 工作集反映进程当前局部性需求，系统应尽量给进程分配足够页框容纳其工作集。
  \item 若分配页框少于工作集需要，缺页率会显著上升。
\end{points}
\plain{工作集不是“这个进程所有页面”，而是“它最近这段时间真正活跃、离不开的页面集合”。}
\exam{老师特别强调了工作集的作用和意义，答题时要点出它是利用局部性来估计当前页面需求，从而降低缺页率、避免抖动。}


\concept{工作集算法}
\begin{points}
  \item 系统定期检查进程在窗口长度 $\Delta$ 内访问过的页面，得到当前工作集。
  \item 若某进程工作集装不下现有驻留集，就应为其增配页框；若系统无空闲页框，可暂停部分进程以释放内存。
  \item 核心思想是让每个进程的当前工作集尽量驻留内存，避免刚换出的页又立刻被访问。
\end{points}
\plain{工作集模型偏“描述当前需要多少页”，工作集算法偏“按这个估计结果去动态分配页框”。}
\exam{名词解释或简答题里，要点出时间窗口 $\Delta$、按工作集调整驻留集、必要时降低多道程序度。}


\concept{抖动 thrashing}
\begin{points}
  \item 抖动指系统花费大量时间进行页面换入换出，真正执行用户程序的时间很少。
  \item 现象：CPU 利用率低，磁盘交换区利用率很高，缺页频繁。
  \item 原因：多道程序度过高或每个进程分得页框不足，工作集无法驻留内存。
  \item 解决：降低多道程序度、增加页框、采用工作集策略、改善置换算法。
\end{points}
\plain{抖动不是简单的“缺页多一点”，而是系统大部分时间都耗在换页上，几乎不干正事。}
\exam{要会和 Belady 异常区分：Belady 异常是某些算法下页框增加反而缺页增多；抖动是整体页框紧张导致系统频繁换页。}


\concept{局部抖动与全局抖动}
\begin{points}
  \item 局部抖动：某个进程分到的页框太少，主要是该进程频繁缺页、反复换页。
  \item 全局抖动：系统中许多进程同时页框不足，彼此争抢内存，导致整个系统频繁换页。
  \item 局部问题若继续提高多道程序度或采用不当的全局置换，可能演变成全局抖动。
\end{points}
\plain{局部抖动是“一个进程转不动”，全局抖动是“整个系统都在疯狂换页”。}
\exam{判断题常考“抖动是单个进程现象还是系统现象”，作答时应区分局部和全局两个层次。}


\concept{利用缺页率判断抖动}
\begin{points}
  \item 缺页率持续偏高，说明进程当前驻留集可能小于其工作集需求。
  \item 可设置缺页率上、下界：高于上界时增配页框或暂停部分进程；低于下界时可适当回收页框。
  \item 若许多进程缺页率同时过高，通常表明系统已接近或进入全局抖动。
\end{points}
\plain{缺页率可以看成内存压力表；不是看某一次缺页，而是看一段时间内是不是一直高。}
\exam{简答题容易问“如何发现抖动”，答题时要写出“监视缺页率”而不只是笼统地说“看系统变慢”。}


\concept{影响缺页率的因素}
\begin{points}
  \item 页面大小：过小会增加页表和 I/O 次数，过大会增加内部碎片和无用调入。
  \item 分配页框数：一般页框越多缺页率越低，但 FIFO 可能有 Belady 异常。
  \item 页面置换算法：OPT 最优，LRU 通常较好，FIFO 简单但可能异常。
  \item 程序局部性：局部性越好，缺页率越低。
  \item 多道程序度：过高会导致每个进程页框不足，引发抖动。
\end{points}
\plain{这一节适合按“页面本身大小、给进程多少页框、怎么置换、程序局部性好不好、系统里进程是不是太多”五个维度来记。}
\exam{主要是概念理解题，要求能解释“为什么这个因素会让缺页率升高或降低”，而不是只会列名词。}
